\counterwithout{section}{chapter}
\setcounter{tocdepth}{1}
\setcounter{secnumdepth}{2}
\hypersetup{bookmarksdepth=2,pdftitle={Fukaya 范畴与 Picard–Lefschetz 理论},pdfauthor={Paul Seidel}}
\renewcommand{\thechapter}{\Roman{chapter}}
\renewcommand{\thesection}{\arabic{section}}
\renewcommand{\thesubsection}{(\arabic{section}\alph{subsection})}
\counterwithout{equation}{chapter}
\counterwithin{equation}{section}
\renewcommand{\theequation}{\thesection.\arabic{equation}}
\makeatletter
% Place page destinations in the shipout foreground so figures cannot discard them.
\let\Hy@EveryPageAnchor\relax
\AddToHook{shipout/foreground}{%
  \ifHy@pageanchor
    \xdef\@currentHpage{page.\thepage}%
    \put(0,-18){\hyper@anchorstart{page.\thepage}\hyper@anchorend}%
  \fi
}
\renewcommand{\theELEGANT@samecnt}{\thesection.\arabic{ELEGANT@samecnt}}
\AtBeginEnvironment{note}{\refstepcounter{ELEGANT@samecnt}}
\renewcommand{\notename}{笔记~\theELEGANT@samecnt}
\AtBeginEnvironment{assumption}{\refstepcounter{ELEGANT@samecnt}}
\renewcommand{\assumptionname}{假设~\theELEGANT@samecnt}
\AtBeginEnvironment{example}{\stepcounter{ELEGANT@samecnt}}
\AtBeginEnvironment{exercise}{\stepcounter{ELEGANT@samecnt}}
\AtBeginEnvironment{problem}{\stepcounter{ELEGANT@samecnt}}
\renewcommand{\theexam}{\theELEGANT@samecnt}
\renewcommand{\theexer}{\theELEGANT@samecnt}
\renewcommand{\theprob}{\theELEGANT@samecnt}
\newcommand{\StatementNumber}[1]{\setcounter{ELEGANT@samecnt}{\numexpr#1-1\relax}}
\makeatother
\newcommand{\SourcePage}[1]{\phantomsection\label{source:#1}\ignorespaces}
\newcommand{\sr}[2]{\hyperref[#1]{#2}}
\newcommand{\secref}[1]{\hyperref[sec:#1]{小节~(#1)}}
\newcommand{\stref}[2]{\hyperref[st:#2]{#1~#2}}
\newcommand{\eref}[1]{\hyperref[eq:#1]{(#1)}}
\newcommand{\bcite}[1]{[\hyperlink{bib:#1}{#1}]}
\newcounter{citetag}
\ExplSyntaxOn
\cs_new_protected:Npn \book_cite:n #1
  {
    \stepcounter{citetag}
    \phantomsection\label{cite:\arabic{citetag}}
    \clist_if_exist:cF {g_book_bib_#1_clist}{\clist_new:c {g_book_bib_#1_clist}}
    \clist_gput_right:cx {g_book_bib_#1_clist}{cite:\arabic{citetag}}
    [\hyperlink{bib:#1}{#1}]
  }
\cs_new_protected:Npn \book_backref:n #1
  {
    \group_begin:
    \seq_clear:N \l_tmpa_seq
    \clist_if_exist:cT {g_book_bib_#1_clist}
      {
        \clist_map_inline:cn {g_book_bib_#1_clist}
          {
            \tl_set:Nx \l_tmpa_tl {\getpagerefnumber{##1}}
            \exp_args:NNV \seq_if_in:NnF \l_tmpa_seq \l_tmpa_tl
              {
                \seq_put_right:NV \l_tmpa_seq \l_tmpa_tl
                \allowbreak\enspace\hyperref[##1]{\pageref*{##1}}
              }
          }
      }
    \group_end:
  }
\renewcommand{\bcite}[1]{\book_cite:n{#1}}
\newcommand{\BibBackref}[1]{{\footnotesize\book_backref:n{#1}}}
\ExplSyntaxOff
\newcommand{\Ai}{A_\infty}
\newcommand{\A}{\mathcal A}
\newcommand{\B}{\mathcal B}
\newcommand{\C}{\mathcal C}
\newcommand{\F}{\mathcal F}
\newcommand{\G}{\mathcal G}
\newcommand{\M}{\mathcal M}
\newcommand{\N}{\mathcal N}
\newcommand{\Y}{\mathcal Y}
\newcommand{\LG}{\mathcal L}
\newcommand{\RG}{\mathcal R}
\newcommand{\sgnmark}{\maltese}
\newcommand{\Q}{\mathcal Q}
\newcommand{\K}{\mathbb K}
\newcommand{\Z}{\mathbb Z}
\newcommand{\R}{\mathbb R}
\newcommand{\CC}{\mathbb C}
\DeclareMathOperator{\Ob}{Ob}
\DeclareMathOperator{\Hom}{Hom}
\DeclareMathOperator{\Tw}{Tw}
\DeclareMathOperator{\Cone}{Cone}
\DeclareMathOperator{\Id}{Id}
\DeclareMathOperator{\id}{id}
\DeclareMathOperator{\im}{im}
\DeclareMathOperator{\HH}{HH}
\DeclareMathOperator{\End}{End}
\newcommand{\nufun}{\operatorname{nu\text{-}fun}}
\newcommand{\modcat}{\operatorname{mod}}
\newcommand{\Ch}{\operatorname{Ch}}
\newcommand{\dg}{\mathrm{dg}}
\newcommand{\CP}{\mathbb{CP}}
\newcommand{\RP}{\mathbb{RP}}
\newcommand{\br}{\sharp}
\newcommand{\eps}{\epsilon}
\newcommand{\ACone}{\mathcal C\!\mathit{one}}
\newcommand{\ModShift}{\mathit{SS}}
\setCJKfamilyfont{zhkai}[BoldFont=SimHei,ItalicFont=KaiTi,BoldItalicFont=SimHei]{KaiTi}
\setCJKfamilyfont{zhfs}[BoldFont=SimHei,ItalicFont=FangSong,BoldItalicFont=SimHei]{FangSong}
\allowdisplaybreaks[2]
\emergencystretch=2em
\usepackage{float}
\newcommand{\SourceDiagram}[2]{\begin{figure}[H]\centering\includegraphics[width=#2\paperwidth]{images/#1.pdf}\end{figure}}
\newcommand{\SourceEquationDiagram}[3]{\begin{figure}[H]\begin{equation}\label{eq:#3}\includegraphics[width=#2\paperwidth]{images/#1.pdf}\end{equation}\end{figure}}

\counterwithout{figure}{chapter}
\counterwithin{figure}{section}
\renewcommand{\thefigure}{\thesection.\arabic{figure}}
\newcommand{\SourceFigure}[3]{\begin{figure}[H]\centering\includegraphics[width=#2\paperwidth]{images/#1.pdf}\caption{}\label{fig:#3}\end{figure}}

\newcommand{\Efamily}{\widetilde{\mathcal E}}
\newcommand{\Sfamily}{\widetilde{\mathcal S}}

% ---- 译者注体例（2026-11 校订版）----
\newif\ifshownotes
\shownotestrue
\newcommand{\znote}[2]{\ifshownotes\latexendnote{\textbf{【原书 p.~#1】}\ #2}\fi}
\renewcommand{\notesname}{译者注}
\renewcommand{\enoteformat}{%
  \rightskip=0pt \leftskip=3.0em \parindent=-3.0em \parskip=4pt plus 1pt
  \leavevmode\makebox[3.0em][l]{\normalfont\bfseries\theenmark.}\ignorespaces}
\renewcommand{\enoteheading}{%
  \cleardoublepage
  \chapter*{译者注}%
  \markboth{译者注}{译者注}%
  \addcontentsline{toc}{chapter}{译者注}%
  \noindent\small 本附录汇集全书译者注。方括号内为原书页码；每条注末给出依据或来源。}%
